%%%PAGE 209 rot=0 legibility=good summary=横置页：洛伦兹力下与变速曲线下R、V关系，平抛轨迹处曲率半径与法向力公式；化学平衡（连锁反应、K总）笔记
%%%BLOCK id=209-01 ch=04 sec=曲线运动·曲率半径与法向力 kind=method alt=11 cont=none y=0.10-0.56
洛伦兹力下 $R\uparrow$ $V\uparrow$

变速曲线下 $R\uparrow$ $V\downarrow$\qquad $F_n=\dfrac{mV^2}{\rho}$
%FIG p209_a 0.04 0.29 0.26 0.50
\notefig[0.3]{figures/p209_a.png}
$F_n=mg\cos\theta$
\[
\begin{cases}
mg\cos\theta=\dfrac{mV^2}{\rho}\\[4pt]
mgy=\dfrac12mV^2\\[4pt]
\tan\theta=2\tan\alpha=\dfrac{2y}{x}
\end{cases}
\qquad
y=\dfrac{g^2x^2}{2V_0^2}
\]
% CHECK: 右侧 y 的表达式分子写作 g²x²，与平抛轨迹 y=gx²/(2v0²) 不一致，请核对
%%%END
%%%BLOCK id=209-02 ch=99 sec=化学·化学平衡与平衡常数 kind=note alt=none cont=none y=0.62-0.93
左=右的反应 有连锁反应时 $P\uparrow$ 反应平衡会移动

$K_{\text{总}}=\dfrac{K_1}{K_2}$\qquad $K_{\text{总}}=K_1\cdot K_2$

\hspace{4.2cm}$K_1=\dfrac{K_{\text{总}}}{K_2}$
%FIG p209_b 0.46 0.71 0.66 0.93
\notefig[0.25]{figures/p209_b.png}
%%%END
%%%PAGEEND 209
